\documentclass[11pt,a4paper]{article}

\usepackage[utf8]{inputenc}
\usepackage[T1]{fontenc}
\usepackage[french]{babel}
\usepackage{lmodern}
\usepackage{microtype}
\usepackage{amsmath,amssymb}
\usepackage{siunitx}
\usepackage{geometry}
\usepackage{booktabs}
\usepackage{enumitem}
\usepackage{xcolor}
\usepackage{hyperref}

\geometry{margin=2.4cm}

\hypersetup{
    colorlinks=true,
    linkcolor=blue!50!black,
    urlcolor=blue!50!black,
    citecolor=blue!50!black
}

\title{Principe de l'algorithme de stitching de courbes SAS}
\author{}
\date{}

\begin{document}

\maketitle

\section{Objectif général}

Dans une expérience de diffusion aux petits angles, une même mesure peut être
obtenue à partir de plusieurs configurations instrumentales et parfois de
plusieurs détecteurs pour une même configuration. Chaque courbe fournit une
mesure de l'intensité diffusée sur une gamme de vecteur de diffusion $Q$ :

\[
    I(Q) \pm \Delta I(Q),
\]

avec éventuellement une estimation de la résolution expérimentale :

\[
    \Delta Q(Q).
\]

L'objectif du stitching est de construire une seule courbe finale continue,
cohérente et exploitable à partir de plusieurs courbes partielles. L'algorithme
ne doit cependant pas utiliser aveuglément toutes les courbes disponibles. Une
courbe de grand angle, par exemple, peut être redondante si elle couvre une
gamme de $Q$ déjà mieux mesurée par une autre configuration.

Le principe général est donc le suivant :

\[
\boxed{
    \text{sélectionner les segments utiles, les raccorder, puis concaténer}
}
\]

\section{Notion de segment}

L'algorithme ne raisonne pas seulement en termes de configuration instrumentale.
Il considère chaque courbe détecteur comme un \emph{segment candidat}. Un
segment correspond donc typiquement à une combinaison :

\[
    \text{configuration} + \text{détecteur}.
\]

Par exemple :

\[
\begin{array}{lll}
    \texttt{c9d0}   & \leftrightarrow & \texttt{config\_9},\ \texttt{detector0}, \\
    \texttt{c10d0}  & \leftrightarrow & \texttt{config\_10},\ \texttt{detector0}, \\
    \texttt{c10d2}  & \leftrightarrow & \texttt{config\_10},\ \texttt{detector2}.
\end{array}
\]

Chaque segment contient quatre colonnes :

\[
    Q,\quad I(Q),\quad \Delta I(Q),\quad \Delta Q(Q).
\]

Pour chaque segment, on calcule aussi des informations globales :

\[
    Q_{\min},\quad Q_{\max},\quad N_{\mathrm{points}},
\]

ainsi que des indicateurs de qualité, par exemple :

\[
    \mathrm{median}\left(\frac{\Delta I}{I}\right),
    \qquad
    \mathrm{median}\left(\frac{\Delta Q}{Q}\right).
\]

\section{Pourquoi ne pas concaténer toutes les courbes ?}

Lorsque plusieurs détecteurs sont disponibles, il arrive que certains segments :

\begin{itemize}[leftmargin=1.5em]
    \item recouvrent presque entièrement une zone déjà couverte par un autre
    segment ;
    \item aient une incertitude relative plus grande ;
    \item présentent une résolution moins favorable ;
    \item introduisent une discontinuité dans la courbe finale ;
    \item n'apportent pas de nouvelle gamme utile en $Q$.
\end{itemize}

Dans ce cas, inclure systématiquement tous les segments peut dégrader la courbe
finale. L'algorithme doit donc être capable de rejeter un segment même si celui-ci
est techniquement présent dans les fichiers d'entrée.

Le critère fondamental est :

\[
\boxed{
    \text{un segment est utile seulement s'il apporte une information nouvelle
    et fiable}
}
\]

\section{Raccord entre deux segments}

Considérons deux segments $A$ et $B$ qui se recouvrent sur une gamme de $Q$.
La question est de savoir si le segment $B$ peut être raccordé au segment $A$
par un simple facteur multiplicatif.

On cherche donc un facteur $s$ tel que :

\[
    I_A(Q) \simeq s\,I_B(Q).
\]

Cette convention est importante : le facteur $s$ est le facteur qui transforme
le segment $B$ vers l'échelle du segment $A$.

\subsection{Travail dans l'espace logarithmique}

Plutôt que d'ajuster directement le rapport $I_A/I_B$, on travaille sur le
log-ratio :

\[
    r(Q) = \ln I_A(Q) - \ln I_B(Q).
\]

Si le raccord est purement multiplicatif, alors :

\[
    r(Q) \simeq \ln s.
\]

Autrement dit, dans la zone de recouvrement, le log-ratio doit former un
plateau. Le problème du raccord devient donc un problème de détection de plateau.

\subsection{Incertitude sur le log-ratio}

Si les incertitudes sur les deux intensités sont indépendantes, l'incertitude
sur le log-ratio est approximativement :

\[
    \sigma_r^2(Q)
    =
    \left(\frac{\Delta I_A(Q)}{I_A(Q)}\right)^2
    +
    \left(\frac{\Delta I_B(Q)}{I_B(Q)}\right)^2.
\]

Cette quantité sert à pondérer les points dans la zone de recouvrement.

\subsection{Estimation du facteur de raccord}

Dans une fenêtre de recouvrement donnée $\Omega$, on ajuste une constante :

\[
    r(Q) \simeq \ln s.
\]

Le meilleur estimateur pondéré est :

\[
    \ln s =
    \frac{\sum_{i\in\Omega} w_i r_i}
         {\sum_{i\in\Omega} w_i},
    \qquad
    w_i = \frac{1}{\sigma_{r,i}^2}.
\]

L'incertitude associée est :

\[
    \sigma_{\ln s}
    =
    \left(
        \sum_{i\in\Omega} w_i
    \right)^{-1/2}.
\]

Le facteur multiplicatif est ensuite :

\[
    s = \exp(\ln s),
\]

et son incertitude est approximativement :

\[
    \sigma_s \simeq s\,\sigma_{\ln s}.
\]

\section{Recherche automatique d'une bonne fenêtre de recouvrement}

Le recouvrement géométrique complet entre deux segments n'est pas toujours la
meilleure zone pour calculer le facteur. Les bords de détecteur, la résolution,
le bruit ou des défauts de normalisation peuvent rendre une partie de l'overlap
moins fiable.

L'algorithme teste donc plusieurs fenêtres candidates dans la zone de
recouvrement. Pour chaque fenêtre, il calcule :

\begin{itemize}[leftmargin=1.5em]
    \item le facteur de raccord $s$ ;
    \item le $\chi^2$ réduit du fit par une constante ;
    \item la pente résiduelle du log-ratio ;
    \item la largeur logarithmique de la fenêtre ;
    \item un indicateur de résolution relative.
\end{itemize}

\subsection{Test de platitude}

Un raccord acceptable doit présenter un log-ratio à peu près constant.
Pour vérifier cela, on ajuste aussi :

\[
    r(Q) = a + b\ln Q.
\]

La pente $b$ doit être compatible avec zéro. On utilise donc le score :

\[
    z_b = \left|\frac{b}{\sigma_b}\right|.
\]

Une valeur petite de $z_b$ indique que la pente n'est pas significative. Une
valeur grande indique que le ratio dérive avec $Q$, donc que la fenêtre ne
correspond pas à un vrai plateau.

\subsection{Rôle du $\chi^2$}

Le $\chi^2$ réduit mesure la compatibilité du log-ratio avec une constante :

\[
    \chi^2_\nu
    =
    \frac{1}{N-1}
    \sum_{i\in\Omega}
    \frac{(r_i-\ln s)^2}{\sigma_{r,i}^2}.
\]

Un $\chi^2_\nu$ trop élevé indique que les fluctuations autour du plateau sont
plus grandes que ce qui est attendu à partir des incertitudes.

Cependant, le $\chi^2$ seul n'est pas suffisant : une dérive lente et régulière
du ratio peut parfois donner un $\chi^2$ acceptable si les incertitudes sont
grandes. C'est pourquoi l'algorithme utilise aussi le critère de pente.

\subsection{Indicateur de résolution}

La résolution n'est pas utilisée pour corriger explicitement les courbes.
L'algorithme ne déconvolue pas les données et ne dégrade pas une courbe pour
imiter la résolution d'une autre.

En revanche, la résolution peut être utilisée comme critère de qualité
secondaire. Pour une paire de segments $A$ et $B$, on définit :

\[
    \rho_\Omega
    =
    \mathrm{median}_{Q\in\Omega}
    \left[
        \sqrt{
            \left(\frac{\Delta Q_A}{Q}\right)^2
            +
            \left(\frac{\Delta Q_B}{Q}\right)^2
        }
    \right].
\]

Une fenêtre où $\rho_\Omega$ est plus faible est préférée, toutes choses égales
par ailleurs.

\subsection{Score d'une fenêtre}

Chaque fenêtre candidate reçoit un score de la forme :

\[
    S_\Omega =
    \chi^2_{\nu,\Omega}
    +
    \lambda_b
    \left(\frac{b_\Omega}{\sigma_{b,\Omega}}\right)^2
    +
    \lambda_w
    \frac{\Delta\ln Q_{\mathrm{total}}}
         {\Delta\ln Q_\Omega}
    +
    \lambda_Q
    \left(\frac{\rho_\Omega}{\rho_{\mathrm{ref}}}\right)^2.
\]

Les différents termes ont les rôles suivants :

\begin{description}[leftmargin=2em]
    \item[$\chi^2_\nu$] pénalise les fenêtres qui ne sont pas compatibles avec
    un plateau.
    \item[Terme de pente] pénalise les ratios qui dérivent avec $Q$.
    \item[Terme de largeur] évite de choisir une fenêtre trop petite simplement
    parce qu'elle donne localement un bon score.
    \item[Terme de résolution] favorise les fenêtres où la résolution relative
    combinée est meilleure.
\end{description}

La meilleure fenêtre est celle qui minimise ce score.

\section{Grille de calcul du ratio}

Deux segments n'ont pas nécessairement les mêmes valeurs de $Q$. Pour calculer
le log-ratio, les intensités doivent donc être évaluées sur une même grille.

L'algorithme peut utiliser une grille commune logarithmique dans la zone de
recouvrement :

\[
    Q_k =
    \exp\left[
        \ln Q_{\min}
        +
        k
        \frac{\ln Q_{\max}-\ln Q_{\min}}{N-1}
    \right].
\]

Les deux courbes sont interpolées sur cette grille commune. Le nombre de points
est choisi de manière conservatrice, typiquement :

\[
    N = \min(N_A, N_B),
\]

où $N_A$ et $N_B$ sont les nombres de points réellement présents dans l'overlap
pour les deux segments. Cela évite de sur-échantillonner artificiellement les
données interpolées.

\section{Sélection des segments utiles}

Une fois que l'on sait tester le raccord entre deux segments, l'étape suivante
consiste à construire une chaîne de segments.

L'algorithme part d'un segment initial, par exemple le segment de plus bas $Q$,
puis cherche quel segment peut prolonger la courbe de la manière la plus
pertinente.

Pour chaque candidat $B$ à ajouter après un segment courant $A$, on calcule :

\begin{itemize}[leftmargin=1.5em]
    \item la qualité du raccord $A\rightarrow B$ ;
    \item la qualité intrinsèque du segment $B$ ;
    \item la nouvelle gamme de $Q$ apportée par $B$.
\end{itemize}

La nouvelle couverture logarithmique est définie par :

\[
    \Delta\ln Q_{\mathrm{new}}
    =
    \max
    \left[
        0,\,
        \ln\left(\frac{Q_{\max,B}}{Q_{\max,A}}\right)
    \right].
\]

Un segment qui ne dépasse pas réellement la gamme déjà couverte n'apporte pas
d'information nouvelle significative.

\subsection{Qualité intrinsèque d'un segment}

Pour un segment seul, on définit un score de qualité simple :

\[
    Q_{\mathrm{seg}}
    =
    \mathrm{median}
    \left[
        \left(\frac{\Delta I}{I}\right)^2
        +
        \lambda_{\mathrm{res}}
        \left(\frac{\Delta Q}{Q}\right)^2
    \right].
\]

Plus ce score est petit, meilleure est la qualité du segment.

\subsection{Coût de transition}

Pour choisir le prochain segment, l'algorithme minimise un coût de transition
du type :

\[
    C(A\rightarrow B)
    =
    S_{A\rightarrow B}
    +
    \lambda_{\mathrm{seg}} Q_{\mathrm{seg},B}
    -
    \lambda_{\mathrm{new}} \Delta\ln Q_{\mathrm{new}}.
\]

Le premier terme mesure la qualité du raccord. Le deuxième pénalise les segments
bruités ou de mauvaise résolution. Le troisième favorise les segments qui
étendent réellement la gamme de $Q$.

Ainsi, un segment est favorisé s'il :

\begin{enumerate}[leftmargin=1.5em]
    \item se raccorde proprement au segment courant ;
    \item possède une bonne qualité intrinsèque ;
    \item apporte une nouvelle gamme de $Q$.
\end{enumerate}

\section{Critères d'acceptation et de rejet}

Une transition entre deux segments est acceptée seulement si elle satisfait des
critères minimaux :

\[
    \chi^2_\nu < \chi^2_{\nu,\max},
\]

\[
    \left|\frac{b}{\sigma_b}\right| < z_{b,\max},
\]

\[
    \Delta\ln Q_{\mathrm{new}} > \Delta\ln Q_{\min}.
\]

Si l'une de ces conditions échoue, la transition est rejetée.

Les principaux statuts de rejet sont :

\begin{description}[leftmargin=2em]
    \item[\texttt{rejected\_bad\_overlap}] le ratio n'est pas suffisamment plat
    ou le $\chi^2$ est trop élevé ;
    \item[\texttt{rejected\_no\_new\_q\_coverage}] le segment n'apporte pas
    assez de nouvelle gamme en $Q$ ;
    \item[\texttt{rejected\_redundant\_or\_no\_new\_q\_coverage}] le segment
    est redondant par rapport à la chaîne déjà retenue ;
    \item[\texttt{rejected\_no\_valid\_transition}] aucun raccord acceptable
    n'a été trouvé pour intégrer ce segment.
\end{description}

Ce mécanisme est essentiel pour éviter d'ajouter automatiquement des détecteurs
grand angle qui n'améliorent pas la courbe finale.

\section{Construction de la courbe finale}

Une fois la chaîne de segments retenue, chaque segment reçoit un facteur global.
Si le premier segment a un facteur global $G_0 = 1$, et si la transition
$A\rightarrow B$ donne un facteur $s_{B\rightarrow A}$, alors :

\[
    G_B = G_A\,s_{B\rightarrow A}.
\]

La courbe corrigée du segment $B$ est donc :

\[
    I_{B,\mathrm{corr}}(Q) = G_B I_B(Q),
\]

et son incertitude :

\[
    \Delta I_{B,\mathrm{corr}}(Q) = G_B \Delta I_B(Q).
\]

La résolution $\Delta Q$ n'est pas multipliée par ce facteur, car il s'agit
d'une incertitude sur l'axe $Q$, pas sur l'intensité :

\[
    \Delta Q_{B,\mathrm{corr}}(Q) = \Delta Q_B(Q).
\]

\subsection{Gestion des zones de recouvrement}

Dans les zones de raccord, l'algorithme ne garde pas nécessairement tout
l'overlap. Il conserve généralement une petite région autour de la fenêtre de
fit pour éviter une coupure trop brutale, mais sans dupliquer excessivement les
points.

Si la fenêtre de fit est :

\[
    [Q_{\mathrm{fit,min}}, Q_{\mathrm{fit,max}}],
\]

on peut conserver une fraction centrale de cette fenêtre en échelle
logarithmique. Cela permet de garder une transition visuellement et
numériquement contrôlée.

\section{Choix de l'échelle de référence}

Par défaut, la courbe finale est exprimée dans l'échelle du premier segment
retenu. Cependant, il est souvent préférable d'exprimer la courbe finale dans
l'échelle d'une configuration de référence choisie par l'utilisateur.

Supposons que les facteurs globaux obtenus soient :

\[
    G_0,\quad G_1,\quad G_2,\ldots
\]

et que le segment de référence ait le facteur :

\[
    G_{\mathrm{ref}}.
\]

On renormalise alors tous les facteurs :

\[
    G_i^{\mathrm{new}}
    =
    \frac{G_i}{G_{\mathrm{ref}}}.
\]

Ainsi, le segment de référence a bien :

\[
    G_{\mathrm{ref}}^{\mathrm{new}} = 1.
\]

Cette opération ne change pas la sélection des segments ni la qualité des
raccords. Elle change seulement l'échelle absolue de la courbe finale.

\section{Résumé de l'algorithme}

L'algorithme complet peut être résumé ainsi :

\begin{enumerate}[leftmargin=1.5em]
    \item Charger toutes les courbes disponibles.
    \item Transformer chaque courbe détecteur en segment candidat.
    \item Calculer les indicateurs de qualité de chaque segment.
    \item Tester les raccords possibles entre paires de segments.
    \item Pour chaque raccord, chercher la meilleure fenêtre de plateau du
    log-ratio.
    \item Rejeter les transitions qui ne sont pas assez plates, trop bruitées
    ou sans nouvelle couverture en $Q$.
    \item Construire progressivement une chaîne de segments utiles.
    \item Rejeter les segments redondants ou non raccordables.
    \item Calculer les facteurs globaux de chaque segment retenu.
    \item Renormaliser éventuellement la courbe sur une configuration de
    référence.
    \item Concaténer les segments retenus pour produire la courbe finale.
\end{enumerate}

\section{Points importants à retenir}

\begin{itemize}[leftmargin=1.5em]
    \item Le raccord est basé sur un facteur multiplicatif entre segments.
    \item Ce facteur est estimé à partir d'un plateau du log-ratio.
    \item Le $\chi^2$ seul ne suffit pas ; il faut aussi contrôler la pente du
    log-ratio.
    \item La résolution $\Delta Q$ n'est pas utilisée pour corriger explicitement
    les courbes, mais elle peut aider à choisir les meilleures fenêtres.
    \item Tous les détecteurs ne sont pas nécessairement utilisés.
    \item Un segment peut être rejeté simplement parce qu'il est redondant.
    \item La configuration de référence fixe l'échelle finale, mais ne change pas
    la sélection des segments.
\end{itemize}

\section{Interprétation pratique}

Une bonne courbe finale n'est pas celle qui contient le plus grand nombre de
points, mais celle qui utilise les points les plus fiables et les plus utiles.
L'algorithme cherche donc un compromis entre continuité, qualité statistique,
résolution et couverture en $Q$.

Dans ce cadre, le stitching n'est pas seulement une opération de collage. C'est
une étape de sélection raisonnée de l'information expérimentale.

\[
\boxed{
    \text{stitching} =
    \text{raccord multiplicatif}
    +
    \text{contrôle de qualité}
    +
    \text{sélection des segments utiles}
}
\]

\end{document}
